\def\dato#1{\ifcsname dato:#1\endcsname\csname dato:#1\endcsname\else\errmessage{dato desconocido: #1}\fi}
\expandafter\def\csname dato:PT0.prev\endcsname{20}%
\expandafter\def\csname dato:PT0.medido\endcsname{32}%
\expandafter\def\csname dato:PT0.estimado\endcsname{15}%
\expandafter\def\csname dato:PT0.pendiente\endcsname{0}%
\expandafter\def\csname dato:PT0.registrado\endcsname{47}%
\expandafter\def\csname dato:PT0.real\endcsname{47}%
\expandafter\def\csname dato:PT0.desv\endcsname{+27}%
\expandafter\def\csname dato:PT0.desvpct\endcsname{+135,0}%
\expandafter\def\csname dato:PT1.prev\endcsname{30}%
\expandafter\def\csname dato:PT1.medido\endcsname{3}%
\expandafter\def\csname dato:PT1.estimado\endcsname{50}%
\expandafter\def\csname dato:PT1.pendiente\endcsname{0}%
\expandafter\def\csname dato:PT1.registrado\endcsname{53}%
\expandafter\def\csname dato:PT1.real\endcsname{53}%
\expandafter\def\csname dato:PT1.desv\endcsname{+23}%
\expandafter\def\csname dato:PT1.desvpct\endcsname{+76,7}%
\expandafter\def\csname dato:PT2.prev\endcsname{25}%
\expandafter\def\csname dato:PT2.medido\endcsname{1}%
\expandafter\def\csname dato:PT2.estimado\endcsname{4}%
\expandafter\def\csname dato:PT2.pendiente\endcsname{0}%
\expandafter\def\csname dato:PT2.registrado\endcsname{5}%
\expandafter\def\csname dato:PT2.real\endcsname{5}%
\expandafter\def\csname dato:PT2.desv\endcsname{$-$20}%
\expandafter\def\csname dato:PT2.desvpct\endcsname{$-$80,0}%
\expandafter\def\csname dato:PT3.prev\endcsname{90}%
\expandafter\def\csname dato:PT3.medido\endcsname{29}%
\expandafter\def\csname dato:PT3.estimado\endcsname{9}%
\expandafter\def\csname dato:PT3.pendiente\endcsname{0}%
\expandafter\def\csname dato:PT3.registrado\endcsname{38}%
\expandafter\def\csname dato:PT3.real\endcsname{38}%
\expandafter\def\csname dato:PT3.desv\endcsname{$-$52}%
\expandafter\def\csname dato:PT3.desvpct\endcsname{$-$57,8}%
\expandafter\def\csname dato:PT4.prev\endcsname{30}%
\expandafter\def\csname dato:PT4.medido\endcsname{11}%
\expandafter\def\csname dato:PT4.estimado\endcsname{3}%
\expandafter\def\csname dato:PT4.pendiente\endcsname{0}%
\expandafter\def\csname dato:PT4.registrado\endcsname{14}%
\expandafter\def\csname dato:PT4.real\endcsname{14}%
\expandafter\def\csname dato:PT4.desv\endcsname{$-$16}%
\expandafter\def\csname dato:PT4.desvpct\endcsname{$-$53,3}%
\expandafter\def\csname dato:PT5.prev\endcsname{25}%
\expandafter\def\csname dato:PT5.medido\endcsname{1}%
\expandafter\def\csname dato:PT5.estimado\endcsname{1}%
\expandafter\def\csname dato:PT5.pendiente\endcsname{0}%
\expandafter\def\csname dato:PT5.registrado\endcsname{2}%
\expandafter\def\csname dato:PT5.real\endcsname{2}%
\expandafter\def\csname dato:PT5.desv\endcsname{$-$23}%
\expandafter\def\csname dato:PT5.desvpct\endcsname{$-$92,0}%
\expandafter\def\csname dato:PT6.prev\endcsname{80}%
\expandafter\def\csname dato:PT6.medido\endcsname{32}%
\expandafter\def\csname dato:PT6.estimado\endcsname{8}%
\expandafter\def\csname dato:PT6.pendiente\endcsname{15}%
\expandafter\def\csname dato:PT6.registrado\endcsname{40}%
\expandafter\def\csname dato:PT6.real\endcsname{55}%
\expandafter\def\csname dato:PT6.desv\endcsname{$-$25}%
\expandafter\def\csname dato:PT6.desvpct\endcsname{$-$31,3}%
\expandafter\def\csname dato:total.prev\endcsname{300}%
\expandafter\def\csname dato:total.medido\endcsname{109}%
\expandafter\def\csname dato:total.estimado\endcsname{90}%
\expandafter\def\csname dato:total.registrado\endcsname{199}%
\expandafter\def\csname dato:total.pendiente\endcsname{15}%
\expandafter\def\csname dato:total.real\endcsname{214}%
\expandafter\def\csname dato:total.desv\endcsname{$-$86}%
\expandafter\def\csname dato:total.desvpct\endcsname{$-$28,7}%
\expandafter\def\csname dato:grupo.PT3a5.prev\endcsname{145}%
\expandafter\def\csname dato:grupo.PT3a5.real\endcsname{54}%
\expandafter\def\csname dato:grupo.PT3a5.desv\endcsname{$-$91}%
\expandafter\def\csname dato:grupo.PT3a5.desvpct\endcsname{$-$62,8}%
\expandafter\def\csname dato:medido.dias\endcsname{34}%
\expandafter\def\csname dato:medido.diasmedia\endcsname{33}%
\expandafter\def\csname dato:medido.mediadia\endcsname{3,2}%
\expandafter\def\csname dato:medido.desde\endcsname{7 de septiembre de 2026}%
\expandafter\def\csname dato:medido.hasta\endcsname{29 de octubre de 2026}%
\expandafter\def\csname dato:est.kickoff.horas\endcsname{4}%
\expandafter\def\csname dato:est.estudioligero.horas\endcsname{20}%
\expandafter\def\csname dato:est.puestadia.horas\endcsname{30}%
\expandafter\def\csname dato:est.ago02.horas\endcsname{3,2}%
\expandafter\def\csname dato:est.ago03.horas\endcsname{3,2}%
\expandafter\def\csname dato:est.ago04.horas\endcsname{3,2}%
\expandafter\def\csname dato:est.ago05.horas\endcsname{3,2}%
\expandafter\def\csname dato:est.ago10.horas\endcsname{3,2}%
\expandafter\def\csname dato:est.sep01.horas\endcsname{3,2}%
\expandafter\def\csname dato:est.sep03.horas\endcsname{3,2}%
\expandafter\def\csname dato:est.sep06.horas\endcsname{3,2}%
\expandafter\def\csname dato:est.sep07.horas\endcsname{3,2}%
\expandafter\def\csname dato:est.sep08.horas\endcsname{3,2}%
\expandafter\def\csname dato:est.sep09.horas\endcsname{3,2}%
\expandafter\def\csname dato:est.defensa.horas\endcsname{15}%
\expandafter\def\csname dato:est.mediamedida.horas\endcsname{35}%
\expandafter\def\csname dato:medido.hueco\endcsname{45}%
\expandafter\def\csname dato:medido.cola\endcsname{15}%
\expandafter\def\csname dato:est.mediamedida.dias\endcsname{11}%
\expandafter\def\csname dato:rol.gestor.tarifa\endcsname{43,88}%
\expandafter\def\csname dato:rol.gestor.horasprev\endcsname{100}%
\expandafter\def\csname dato:rol.gestor.horasreal\endcsname{102}%
\expandafter\def\csname dato:rol.gestor.costeprev\endcsname{4.388,00}%
\expandafter\def\csname dato:rol.gestor.costereal\endcsname{4.475,76}%
\expandafter\def\csname dato:rol.gestor.horasreg\endcsname{87}%
\expandafter\def\csname dato:rol.gestor.horaspend\endcsname{15}%
\expandafter\def\csname dato:rol.gestor.costereg\endcsname{3.817,56}%
\expandafter\def\csname dato:rol.gestor.costepend\endcsname{658,20}%
\expandafter\def\csname dato:rol.gestor.desvh\endcsname{+2}%
\expandafter\def\csname dato:rol.gestor.desvcoste\endcsname{+87,76}%
\expandafter\def\csname dato:rol.analista.tarifa\endcsname{35,09}%
\expandafter\def\csname dato:rol.analista.horasprev\endcsname{110}%
\expandafter\def\csname dato:rol.analista.horasreal\endcsname{74}%
\expandafter\def\csname dato:rol.analista.costeprev\endcsname{3.859,90}%
\expandafter\def\csname dato:rol.analista.costereal\endcsname{2.596,66}%
\expandafter\def\csname dato:rol.analista.horasreg\endcsname{74}%
\expandafter\def\csname dato:rol.analista.horaspend\endcsname{0}%
\expandafter\def\csname dato:rol.analista.costereg\endcsname{2.596,66}%
\expandafter\def\csname dato:rol.analista.costepend\endcsname{0,00}%
\expandafter\def\csname dato:rol.analista.desvh\endcsname{$-$36}%
\expandafter\def\csname dato:rol.analista.desvcoste\endcsname{$-$1.263,24}%
\expandafter\def\csname dato:rol.desarrollador.tarifa\endcsname{27,73}%
\expandafter\def\csname dato:rol.desarrollador.horasprev\endcsname{90}%
\expandafter\def\csname dato:rol.desarrollador.horasreal\endcsname{38}%
\expandafter\def\csname dato:rol.desarrollador.costeprev\endcsname{2.495,70}%
\expandafter\def\csname dato:rol.desarrollador.costereal\endcsname{1.053,74}%
\expandafter\def\csname dato:rol.desarrollador.horasreg\endcsname{38}%
\expandafter\def\csname dato:rol.desarrollador.horaspend\endcsname{0}%
\expandafter\def\csname dato:rol.desarrollador.costereg\endcsname{1.053,74}%
\expandafter\def\csname dato:rol.desarrollador.costepend\endcsname{0,00}%
\expandafter\def\csname dato:rol.desarrollador.desvh\endcsname{$-$52}%
\expandafter\def\csname dato:rol.desarrollador.desvcoste\endcsname{$-$1.441,96}%
\expandafter\def\csname dato:amort.precio\endcsname{799,00}%
\expandafter\def\csname dato:amort.vida\endcsname{60}%
\expandafter\def\csname dato:amort.meses\endcsname{9}%
\expandafter\def\csname dato:fijo.claude.mensual\endcsname{90,00}%
\expandafter\def\csname dato:fijo.claude.meses\endcsname{1}%
\expandafter\def\csname dato:fijo.claude.total\endcsname{90,00}%
\expandafter\def\csname dato:fijo.github.mensual\endcsname{3,58}%
\expandafter\def\csname dato:fijo.github.meses\endcsname{8}%
\expandafter\def\csname dato:fijo.github.total\endcsname{28,64}%
\expandafter\def\csname dato:fijo.internet.mensual\endcsname{20,00}%
\expandafter\def\csname dato:fijo.internet.meses\endcsname{8}%
\expandafter\def\csname dato:fijo.internet.total\endcsname{160,00}%
\expandafter\def\csname dato:fijo.luz.mensual\endcsname{25,00}%
\expandafter\def\csname dato:fijo.luz.meses\endcsname{8}%
\expandafter\def\csname dato:fijo.luz.total\endcsname{200,00}%
\expandafter\def\csname dato:pres.contingenciapct\endcsname{10}%
\expandafter\def\csname dato:pres.personal.prev\endcsname{10.743,60}%
\expandafter\def\csname dato:pres.personal.real\endcsname{8.126,16}%
\expandafter\def\csname dato:pres.personal.reg\endcsname{7.467,96}%
\expandafter\def\csname dato:pres.personal.pend\endcsname{658,20}%
\expandafter\def\csname dato:pres.personal.desv\endcsname{$-$2.617,44}%
\expandafter\def\csname dato:pres.personal.desvpct\endcsname{$-$24,4}%
\expandafter\def\csname dato:pres.amort\endcsname{119,85}%
\expandafter\def\csname dato:pres.fijos\endcsname{478,64}%
\expandafter\def\csname dato:pres.subtotal.prev\endcsname{11.342,09}%
\expandafter\def\csname dato:pres.subtotal.real\endcsname{8.724,65}%
\expandafter\def\csname dato:pres.contingencia\endcsname{1.134,21}%
\expandafter\def\csname dato:pres.total.prev\endcsname{12.476,30}%
\expandafter\def\csname dato:pres.desv\endcsname{$-$2.617,44}%
\expandafter\def\csname dato:pres.desvpct\endcsname{$-$23,1}%
\expandafter\def\csname dato:pres.total.desv\endcsname{$-$3.751,65}%
\expandafter\def\csname dato:pres.total.desvpct\endcsname{$-$30,1}%

\chapter{Metodología y planificación}
\label{cap:metodologia}

Este capítulo explica cómo se ha hecho el trabajo. Primero presenta la metodología
(Sección~\ref{sec:metodologia:enfoque}). Después recoge la planificación,
el presupuesto y las desviaciones respecto al plan
(Secciones~\ref{sec:metodologia:planificacion}
a~\ref{sec:metodologia:desviaciones}).

\section{Metodología}
\label{sec:metodologia:enfoque}

El trabajo adopta el enfoque de investigación en ciencia del diseño
(\textit{Design Science Research}, DSR) formulado por Hevner et al.
\citep{hevner2004dsr}. Se elige porque el objetivo de la
Sección~\ref{sec:introduccion:objetivos}, un catálogo de Datasheets y un
método de comparación, pide construir artefactos que resuelvan un problema relevante y evaluar con rigor su utilidad, y no
describir un fenómeno. Para aplicar el enfoque se sigue el
proceso DSRM (\textit{Design Science Research Methodology}) de Peffers
et al. \citep{peffers2007dsrm}, que ordena la
investigación en seis actividades y admite iterar entre ellas.

El trabajo tiene cuatro contribuciones (C-01 a C-04). Tres son artefactos de alguno de los cuatro tipos que distinguen Hevner et al.: constructos, modelos, métodos e instanciaciones. El catálogo de Datasheets es una instanciación (C-01). Las extensiones al esquema son un modelo (C-02), construido sobre el ACM y su serialización ACM2YAML, de Gavira Sánchez \citep{gavira2026tfm}. El método de comparación sistemática es un método (C-04). El análisis de límites y \textit{rates} (C-03) no es un artefacto: es lo que se aprende al aplicar las operaciones de PRIME4API al catálogo.

La Tabla~\ref{tab:metodologia:dsrm} recorre las seis actividades del proceso DSRM. Para cada una dice qué se hace en este trabajo, en qué capítulo se recoge y a qué objetivo responde.

Cada Datasheet se construye con un ciclo de cinco fases: recolección de fuentes, borrador, validación, revisión y consolidación. La validación es automática, con el validador de formato, y la revisión examina el borrador campo a campo con las fuentes a la vista (Sección~\ref{sec:propuesta:ciclo}). El ciclo se probó y se ajustó con Postmark, la API piloto, antes de aplicarlo al resto de APIs. Las decisiones de modelado que se resuelven, como qué se considera cuota, se anotan como reglas en el libro de criterios, un fichero de reglas común a todas las Datasheets, y no se vuelven a plantear. Las decisiones de alcance se anotan en un registro, y de él sale la comparación entre lo previsto y lo ocurrido de la Sección~\ref{sec:metodologia:desviaciones}.

El trabajo lo ejecutan agentes de inteligencia artificial basados en Claude, el asistente de Anthropic (Anexo~\ref{anx:uso-ia}), y lo supervisa el autor. Los agentes recogen las fuentes, construyen cada Datasheet y la verifican, en flujos paralelos. El autor diseña el procedimiento, fija los criterios de modelado (como qué fuente prevalece cuando dos se contradicen), resuelve lo que queda abierto y valida el resultado. La verificación la hace un agente distinto del que construye la Datasheet, y reabre cada URL citada. La Datasheet de Postmark, la del piloto, se construyó antes de que se separaran la construcción y la verificación, y se cerró sin informe de verificador. Las otras siete lo tienen. Postmark sí pasó por el validador de formato y por PRIME4API (Tabla~\ref{tab:verificacion}).

\begin{table}[!htbp]
\centering
\small
\setlength{\tabcolsep}{3pt}
\begin{tabular}{|p{0.2\textwidth}|p{0.43\textwidth}|p{0.15\textwidth}|p{0.13\textwidth}|}
\hline
\raggedright \textbf{Actividad} & \raggedright \textbf{En este trabajo} & \raggedright \textbf{Dónde} & \raggedright \textbf{Objetivo} \tabularnewline
\hline
\raggedright A1. Problema y motivación & \raggedright Capacidad contractual dispersa en fuentes heterogéneas, e instrumental validado sobre pocos servicios & \raggedright Caps.~\ref{cap:introduccion} y~\ref{cap:estadoDelArte} & \raggedright --- \tabularnewline
\hline
\raggedright A2. Objetivos de la solución & \raggedright Preguntas RQ1 a RQ4 y objetivos específicos, contra los que se evalúa cada actividad posterior & \raggedright Sec.~\ref{sec:introduccion:objetivos} & \raggedright OBJ-01 a OBJ-06 \tabularnewline
\hline
\raggedright A3. Diseño y desarrollo & \raggedright Criterio de selección, ciclo de construcción por API, trazabilidad de cada valor a su fuente y escalera de validación. Sobre el catálogo terminado, las extensiones del esquema y el catálogo de preguntas de decisión & \raggedright Caps.~\ref{cap:catalogo} y~\ref{cap:evolucion}, Sec.~\ref{sec:comparacion:catalogo} & \raggedright OBJ-01, OBJ-02 y OBJ-03 \tabularnewline
\hline
\raggedright A4. Demostración & \raggedright Las ocho Datasheets se ejecutan con PRIME4API y se preparan como contribución a SPHERE. Sobre ellas se analiza la capacidad y se aplica el método de comparación & \raggedright Tabla~\ref{tab:verificacion}, Sec.~\ref{sec:propuesta:sphere}, Caps.~\ref{cap:limites} y~\ref{cap:comparacion} & \raggedright OBJ-04, OBJ-05 y OBJ-06 \tabularnewline
\hline
\raggedright A5. Evaluación & \raggedright Validador de formato, verificación por uso sobre PRIME4API y reapertura de cada URL citada. Criterio de promoción, que decide qué elementos no soportados (lo que el esquema no expresa o PRIME4API no consume) entran en el modelo, y clase de cada pregunta, que dice con qué se responde o qué falta para responderla & \raggedright Caps.~\ref{cap:catalogo} a~\ref{cap:comparacion} & \raggedright Todos \tabularnewline
\hline
\raggedright A6. Comunicación & \raggedright Esta memoria y el repositorio de datos abierto & \raggedright Cap.~\ref{cap:conclusiones} & \raggedright --- \tabularnewline
\hline
\end{tabular}
\caption[Actividades del proceso DSRM]{Actividades del proceso DSRM \citep{peffers2007dsrm}, con lo que cada una hace en este trabajo, el capítulo que la recoge y el objetivo al que responde.}
\label{tab:metodologia:dsrm}
\end{table}

El límite del enfoque está en la evaluación. El ecosistema no tiene un validador oficial de Datasheets. El validador de formato es de este trabajo: contrasta cada Datasheet con la guía de ACM2YAML v0.3 y comprueba los tipos de cada campo y las referencias entre bloques, como que exista la cuota que cita un plan. La ejecución en PRIME4API comprueba que el motor consume la Datasheet y que sus respuestas coinciden con el cálculo a mano en las preguntas que lo permiten. La utilidad para quien elige un proveedor, que es lo que DSR pide evaluar, no se ha medido, porque no hay un estudio con integradores, los desarrolladores que conectan su aplicación a la API de un proveedor.

\section{Planificación}
\label{sec:metodologia:planificacion}

El Trabajo Fin de Máster tiene 12 créditos ECTS \citep{us2026tfm}. A
25 horas por crédito, el mínimo del Real Decreto 1125/2003
\citep{rd1125de2003}, el esfuerzo previsto es de \dato{total.prev}~horas.
El proyecto empieza el 13 de febrero de 2026. La línea base reparte las
horas entre siete paquetes de trabajo
(Tabla~\ref{tab:metodologia:paquetes}). PT0 y PT6 son transversales, y
PT1 a PT5 siguen el orden del método. Cada paquete tiene un entregable,
el objetivo específico al que responde y un rol, que fija su tarifa en
el presupuesto.

\begin{table}[tbp]
\centering
\small
\setlength{\tabcolsep}{3pt}
\begin{tabular}{|p{0.19\textwidth}|p{0.37\textwidth}|p{0.12\textwidth}|p{0.15\textwidth}|r|}
\hline
\raggedright \textbf{Paquete} & \raggedright \textbf{Entregable} & \raggedright \textbf{Objetivo} & \raggedright \textbf{Rol} & \textbf{Horas} \tabularnewline
\hline
\raggedright PT0. Gestión y seguimiento & \raggedright Registro de decisiones, seguimiento del plan y repositorio versionado & \raggedright --- & \raggedright Gestor de proyecto & \dato{PT0.prev} \tabularnewline
\hline
\raggedright PT1. Estudio previo & \raggedright Estado del arte sobre modelado de pricings y de capacidad de APIs & \raggedright --- & \raggedright Analista de sistemas & \dato{PT1.prev} \tabularnewline
\hline
\raggedright PT2. Selección y piloto & \raggedright Criterio de inclusión, conjunto de APIs y Datasheet piloto & \raggedright OBJ-01 & \raggedright Analista de sistemas & \dato{PT2.prev} \tabularnewline
\hline
\raggedright PT3. Construcción del catálogo & \raggedright Una Datasheet por API, con fuentes, observación de panel y sondeos. Inventario de elementos no soportados por el esquema y propuestas de extensión & \raggedright OBJ-02 y OBJ-03 & \raggedright Desarrollador & \dato{PT3.prev} \tabularnewline
\hline
\raggedright PT4. Límites y \textit{rates} & \raggedright Datasheets ejecutadas en PRIME4API y caracterización de la relación entre límites y \textit{rates} & \raggedright OBJ-04 & \raggedright Analista de sistemas & \dato{PT4.prev} \tabularnewline
\hline
\raggedright PT5. Comparación sistemática & \raggedright Catálogo de preguntas de decisión y procedimiento para responderlas & \raggedright OBJ-05 & \raggedright Analista de sistemas & \dato{PT5.prev} \tabularnewline
\hline
\raggedright PT6. Memoria y defensa & \raggedright Memoria depositada y defensa & \raggedright --- & \raggedright Gestor de proyecto & \dato{PT6.prev} \tabularnewline
\hline
\raggedright \textbf{Total} & & & & \textbf{\dato{total.prev}} \tabularnewline
\hline
\end{tabular}
\caption[Paquetes de trabajo de la línea base]{Paquetes de trabajo de la línea base, con su entregable, el objetivo específico de la Sección~\ref{sec:introduccion:objetivos} al que responde, su rol y las horas previstas.}
\label{tab:metodologia:paquetes}
\end{table}

La línea base distribuye el trabajo entre febrero y septiembre
(Figura~\ref{fig:metodologia:cronograma}): selección y piloto hasta el
15 de abril, catálogo hasta el 15 de julio, límites y comparación hasta
el 31 de julio y memoria hasta el 30 de septiembre. La línea base se
reconstruyó el 29 de septiembre de 2026: el reparto de horas y las
ventanas son una asignación nominal del total, y el primer plan escrito
con fechas es del 2 de agosto. Por eso, las desviaciones de calendario
frente a ella miden la distancia a ese reparto y no el incumplimiento de
un compromiso previo.

El 1 de septiembre se replanifica el trabajo restante con un calendario
semanal. Ese calendario fija el cierre del catálogo el 28 de septiembre,
la versión final de la memoria el 12 de octubre y el depósito el 15 de
octubre. El 30 de septiembre, el calendario oficial de la tercera
convocatoria sitúa el depósito el 2 de noviembre y la defensa entre el 9
y el 13 de noviembre de 2026.

\begin{figure}[tbp]
\centering
\begin{tikzpicture}[x=1.3cm,y=1cm,font=\scriptsize]
\node[anchor=south] at (0.5,0.1) {feb 26};
\node[anchor=south] at (1.5,0.1) {mar 26};
\node[anchor=south] at (2.5,0.1) {abr 26};
\node[anchor=south] at (3.5,0.1) {may 26};
\node[anchor=south] at (4.5,0.1) {jun 26};
\node[anchor=south] at (5.5,0.1) {jul 26};
\node[anchor=south] at (6.5,0.1) {ago 26};
\node[anchor=south] at (7.5,0.1) {sep 26};
\node[anchor=south] at (8.5,0.1) {oct 26};
\node[anchor=south] at (9.5,0.1) {nov 26};
\draw[gray!40] (0,0) -- (0,-7);
\draw[gray!40] (1,0) -- (1,-7);
\draw[gray!40] (2,0) -- (2,-7);
\draw[gray!40] (3,0) -- (3,-7);
\draw[gray!40] (4,0) -- (4,-7);
\draw[gray!40] (5,0) -- (5,-7);
\draw[gray!40] (6,0) -- (6,-7);
\draw[gray!40] (7,0) -- (7,-7);
\draw[gray!40] (8,0) -- (8,-7);
\draw[gray!40] (9,0) -- (9,-7);
\draw[gray!40] (10,0) -- (10,-7);
\node[anchor=east] at (0,-0.5) {PT0};
\draw (0.429,-0.10) rectangle (8.000,-0.45);
\fill[black,opacity=0.350] (0,-0.55) rectangle (1,-0.90);
\fill[black,opacity=0.424] (6,-0.55) rectangle (7,-0.90);
\fill[black,opacity=0.756] (7,-0.55) rectangle (8,-0.90);
\fill[black,opacity=0.641] (8,-0.55) rectangle (9,-0.90);
\node[anchor=east] at (0,-1.5) {PT1};
\draw (0.429,-1.10) rectangle (2.000,-1.45);
\fill[black,opacity=0.349] (0,-1.55) rectangle (1,-1.90);
\fill[black,opacity=0.454] (1,-1.55) rectangle (2,-1.90);
\fill[black,opacity=0.447] (2,-1.55) rectangle (3,-1.90);
\fill[black,opacity=1.000] (3,-1.55) rectangle (4,-1.90);
\fill[black,opacity=0.310] (7,-1.55) rectangle (8,-1.90);
\fill[black,opacity=0.266] (8,-1.55) rectangle (9,-1.90);
\node[anchor=east] at (0,-2.5) {PT2};
\draw (1.000,-2.10) rectangle (2.500,-2.45);
\fill[black,opacity=0.337] (6,-2.55) rectangle (7,-2.90);
\fill[black,opacity=0.270] (7,-2.55) rectangle (8,-2.90);
\fill[black,opacity=0.264] (8,-2.55) rectangle (9,-2.90);
\node[anchor=east] at (0,-3.5) {PT3};
\draw (2.000,-3.10) rectangle (5.484,-3.45);
\fill[black,opacity=0.330] (6,-3.55) rectangle (7,-3.90);
\fill[black,opacity=0.960] (7,-3.55) rectangle (8,-3.90);
\fill[black,opacity=0.417] (8,-3.55) rectangle (9,-3.90);
\node[anchor=east] at (0,-4.5) {PT4};
\draw (3.452,-4.10) rectangle (6.000,-4.45);
\fill[black,opacity=0.263] (6,-4.55) rectangle (7,-4.90);
\fill[black,opacity=0.434] (7,-4.55) rectangle (8,-4.90);
\fill[black,opacity=0.392] (8,-4.55) rectangle (9,-4.90);
\node[anchor=east] at (0,-5.5) {PT5};
\draw (4.467,-5.10) rectangle (6.000,-5.45);
\fill[black,opacity=0.282] (7,-5.55) rectangle (8,-5.90);
\fill[black,opacity=0.264] (8,-5.55) rectangle (9,-5.90);
\node[anchor=east] at (0,-6.5) {PT6};
\draw (1.000,-6.10) rectangle (8.000,-6.45);
\fill[black,opacity=0.297] (6,-6.55) rectangle (7,-6.90);
\fill[black,opacity=0.838] (7,-6.55) rectangle (8,-6.90);
\fill[black,opacity=0.648] (8,-6.55) rectangle (9,-6.90);
\fill[black,opacity=0.575] (9,-6.55) rectangle (10,-6.90);
\end{tikzpicture}

\caption[Cronograma previsto y al cierre]{Cronograma del proyecto. Por cada paquete, la barra sin relleno es su ventana en la línea base, y las celdas rellenas de debajo son las horas de cada mes, más intensas cuantas más horas. Las \dato{total.sinregistro} horas de revisión e iteración posteriores al primer borrador completo no tienen mes asignado y no aparecen.}
\label{fig:metodologia:cronograma}
\end{figure}

\section{Presupuesto}
\label{sec:metodologia:presupuesto}

El presupuesto tiene tres partidas: personal, amortización del equipo y
costes fijos, que son licencias y suministros. Sobre el subtotal previsto se añade una contingencia del
\dato{pres.contingenciapct}~\%. El cálculo se apoya en las siguientes
hipótesis.

\begin{description}
  \item[Horas.] Las previstas son las de la línea base y suman
    \dato{total.prev}. Las de cierre suman \dato{total.real}
    (Tabla~\ref{tab:metodologia:horas}). Las horas de cada rol son la suma
    de las de sus paquetes.
  \item[Tarifas.] Cada rol se valora con la parte de costes directos,
    el 77,46~\%, del precio medio de referencia que fija la Instrucción
    1/2023 de la Agencia Digital de Andalucía \citep{ada2023perfiles}.
    Ese precio es el de una empresa contratista, y la parte de costes
    directos excluye sus gastos generales y su beneficio industrial. La
    cotización a la Seguridad Social ya está incluida.
  \item[Amortización.] El equipo es un portátil Huawei MateBook D~16
    (2024). Se toma como coste de adquisición el precio de partida de la
    gama, \dato{amort.precio}~€ \citep{xataka2024matebook}, con una vida
    útil de \dato{amort.vida}~meses. Se amortiza linealmente durante los
    \dato{amort.meses}~meses que van del inicio a la defensa.
  \item[Licencias.] Se imputan dos: un mes de suscripción a Claude
    Max 5x, por su importe facturado (\dato{fijo.claude.mensual}~€), y
    GitHub Pro, el plan individual de pago de GitHub, donde se aloja el
    repositorio de datos. GitHub Pro se valora a su precio de lista,
    aunque el autor no lo paga por ser estudiante: 4~USD al mes, el precio
    que GitHub fijó en abril de 2020 según su documentación
    \citep{githubfaqplanes}. Al tipo de cambio de referencia del Banco
    Central Europeo \citep{bce2026usd} del 30 de octubre de 2026
    (1~€ = 1,1172~USD), son \dato{fijo.github.mensual}~€ al mes. GitHub
    Pro cuenta los mismos \dato{fijo.github.meses}~meses que los
    suministros: el repositorio de datos se creó en octubre, pero el
    trabajo debió alojarse en GitHub desde el principio. Los proveedores del
    catálogo no suman coste, porque se usaron sus planes gratuitos o de
    prueba, abiertos a cualquiera.
  \item[Suministros.] Internet (\dato{fijo.internet.mensual}~€ al mes)
    y electricidad (\dato{fijo.luz.mensual}~€ al mes) cuentan en los
    \dato{fijo.internet.meses}~meses con horas,
    de febrero a mayo y de agosto a noviembre. PRIME4API y SPHERE se
    ejecutan en local, en el mismo equipo.
  \item[Costes que no dependen del número de horas.] La amortización y los
    costes fijos entran igual en la columna prevista y en la de cierre.
\end{description}

La amortización resulta de
\begin{equation}
  C_{\text{amortización}} = \frac{\text{\dato{amort.precio}}}{\text{\dato{amort.vida}}} \times \text{\dato{amort.meses}} = \text{\dato{pres.amort}~€}
\end{equation}

\begin{table}[!htbp]
\centering
\footnotesize
\setlength{\tabcolsep}{3pt}
\begin{tabular}{|l|r|r|r|r|}
\hline
\textbf{Partida} & \textbf{\shortstack{Horas\\prev./cierre}} & \textbf{Previsto} & \textbf{Al cierre} & \textbf{Desviación} \\
\hline
Gestor de proyecto (\dato{rol.gestor.tarifa}) & \dato{rol.gestor.horasprev} / \dato{rol.gestor.horasreal} & \dato{rol.gestor.costeprev} & \dato{rol.gestor.costereal} & \dato{rol.gestor.desvcoste} \\
Analista de sistemas (\dato{rol.analista.tarifa}) & \dato{rol.analista.horasprev} / \dato{rol.analista.horasreal} & \dato{rol.analista.costeprev} & \dato{rol.analista.costereal} & \dato{rol.analista.desvcoste} \\
Desarrollador (\dato{rol.desarrollador.tarifa}) & \dato{rol.desarrollador.horasprev} / \dato{rol.desarrollador.horasreal} & \dato{rol.desarrollador.costeprev} & \dato{rol.desarrollador.costereal} & \dato{rol.desarrollador.desvcoste} \\
\textit{Personal} & \textit{\dato{total.prev} / \dato{total.real}} & \textit{\dato{pres.personal.prev}} & \textit{\dato{pres.personal.real}} & \textit{\dato{pres.personal.desv}} \\
\hline
Amortización del equipo & & \dato{pres.amort} & \dato{pres.amort} & 0,00 \\
Claude Max 5x (\dato{fijo.claude.meses} mes) & & \dato{fijo.claude.total} & \dato{fijo.claude.total} & 0,00 \\
GitHub Pro (\dato{fijo.github.meses} meses) & & \dato{fijo.github.total} & \dato{fijo.github.total} & 0,00 \\
Internet (\dato{fijo.internet.meses} meses) & & \dato{fijo.internet.total} & \dato{fijo.internet.total} & 0,00 \\
Electricidad (\dato{fijo.luz.meses} meses) & & \dato{fijo.luz.total} & \dato{fijo.luz.total} & 0,00 \\
\hline
Subtotal & & \dato{pres.subtotal.prev} & \dato{pres.subtotal.real} & \dato{pres.desv} \\
Contingencia (\dato{pres.contingenciapct}~\%) & & \dato{pres.contingencia} & --- & \\
\hline
\textbf{Total} & & \textbf{\dato{pres.total.prev}} & \textbf{\dato{pres.subtotal.real}} & \textbf{\dato{pres.total.desv}} \\
\hline
\end{tabular}
\caption[Presupuesto, previsto y al cierre]{Presupuesto previsto y al cierre, en euros sin IVA. Entre paréntesis, la tarifa de cada rol en euros por hora. La columna al cierre valora las horas al cierre y no incluye contingencia.}
\label{tab:metodologia:total}
\end{table}

El subtotal al cierre se desvía \dato{pres.desv}~€ del previsto
(\dato{pres.desvpct}~\%). Como los costes que no dependen del número de horas
son iguales en las dos columnas, toda la diferencia es de personal
(\dato{pres.personal.desvpct}~\%). Por roles, la desviación es de
\dato{rol.gestor.desvh}~horas en el gestor de proyecto,
\dato{rol.analista.desvh} en el analista y \dato{rol.desarrollador.desvh}
en el desarrollador. La desviación del subtotal cabe en la contingencia,
de \dato{pres.contingencia}~€, y el total al cierre se desvía
\dato{pres.total.desv}~€ del total previsto
(\dato{pres.total.desvpct}~\%).

\section{Desviaciones}
\label{sec:metodologia:desviaciones}

La ejecución se compara con la planificación según cuatro criterios por
separado: calendario, horas, entregables y alcance. El resultado de un
criterio no indica el de otro. Un paquete puede cerrar su entregable con menos
horas de las previstas y fuera de su ventana.

\paragraph{Calendario.} La Tabla~\ref{tab:metodologia:hitos} compara
las fechas de los hitos con la línea base y con la replanificación del
1 de septiembre. Respecto a la línea base, el catálogo se completa 68
días más tarde y la memoria, 29 días más tarde. Respecto a la
replanificación, el catálogo se adelanta 7 días, y la memoria y el
depósito se retrasan 17 y 18 días. De las horas con mes asignado, las de febrero a mayo son las de la reunión de arranque
(\dato{est.kickoff.horas}~horas) y las del estudio previo:
\dato{est.estudioligero.horas} de lectura del estado del arte hasta abril
y \dato{est.puestadia.horas} de puesta al día en las líneas de
investigación del grupo en mayo. Junio y julio no tienen horas, y el
resto se concentra entre agosto y octubre
(Figura~\ref{fig:metodologia:cronograma}).

\begin{table}[!htbp]
\centering
\small
\setlength{\tabcolsep}{4pt}
\begin{tabular}{|l|c|c|c|r|r|}
\hline
\textbf{Hito} & \textbf{\shortstack{Línea\\base}} & \textbf{\shortstack{Replanif.\\1 sep}} & \textbf{\shortstack{Real o\\previsto}} & \textbf{\shortstack{Días frente a\\línea base}} & \textbf{\shortstack{Días frente a\\replanif.}} \\
\hline
Catálogo completo & 15 jul & 28 sep & 21 sep & +68 & $-$7 \\
Memoria terminada & 30 sep & 12 oct & 29 oct\textsuperscript{a} & +29 & +17 \\
Depósito & --- & 15 oct & 2 nov\textsuperscript{a} & --- & +18 \\
Defensa & --- & --- & 9--13 nov\textsuperscript{a} & --- & --- \\
\hline
\end{tabular}
\caption[Hitos previstos y reales]{Fechas de los hitos en 2026. La línea base no fija fecha de depósito ni de defensa. \textsuperscript{a}~Fecha prevista.}
\label{tab:metodologia:hitos}
\end{table}

\paragraph{Horas.} La planificación inicial es de \dato{total.prev}~horas,
y la dedicación real, de \dato{total.real}: \dato{total.desv}~horas
(\dato{total.desvpct}~\%, Tabla~\ref{tab:metodologia:horas}). La revisión
y las iteraciones posteriores al primer borrador completo llevaron más
tiempo del previsto. Suben la gestión (PT0), el estudio previo (PT1) y la memoria
(PT6), y bajan los paquetes PT2 a PT5. Las actividades de PT4 y PT5 se
ejecutaron dentro del ciclo de construcción de cada Datasheet y de la
redacción de la memoria, así que la desviación de PT3, PT4 y PT5 por
separado no mide el esfuerzo de cada uno. Juntos pasan de
\dato{grupo.PT3a5.prev} a \dato{grupo.PT3a5.real}~horas
(\dato{grupo.PT3a5.desv}, \dato{grupo.PT3a5.desvpct}~\%).

\begin{table}[!htbp]
\centering
\small
\setlength{\tabcolsep}{4pt}
\begin{tabular}{|l|r|r|r|r|}
\hline
\textbf{Paquete} & \textbf{Previstas} & \textbf{Al cierre} & \textbf{Desv.} & \textbf{Desv. (\%)} \\
\hline
PT0. Gestión & \dato{PT0.prev} & \dato{PT0.real} & \dato{PT0.desv} & \dato{PT0.desvpct} \\
PT1. Estudio previo & \dato{PT1.prev} & \dato{PT1.real} & \dato{PT1.desv} & \dato{PT1.desvpct} \\
PT2. Selección y piloto & \dato{PT2.prev} & \dato{PT2.real} & \dato{PT2.desv} & \dato{PT2.desvpct} \\
PT3. Catálogo & \dato{PT3.prev} & \dato{PT3.real} & \dato{PT3.desv} & \dato{PT3.desvpct} \\
PT4. Límites y \textit{rates} & \dato{PT4.prev} & \dato{PT4.real} & \dato{PT4.desv} & \dato{PT4.desvpct} \\
PT5. Comparación & \dato{PT5.prev} & \dato{PT5.real} & \dato{PT5.desv} & \dato{PT5.desvpct} \\
PT6. Memoria y defensa & \dato{PT6.prev} & \dato{PT6.real} & \dato{PT6.desv} & \dato{PT6.desvpct} \\
\hline
\textbf{Total} & \textbf{\dato{total.prev}} & \textbf{\dato{total.real}} & \textbf{\dato{total.desv}} & \textbf{\dato{total.desvpct}} \\
\hline
\end{tabular}
\caption[Horas previstas y al cierre por paquete]{Horas previstas y al cierre por paquete de trabajo. Las de PT6 incluyen la preparación de la defensa.}
\label{tab:metodologia:horas}
\end{table}

\paragraph{Entregables.} La Tabla~\ref{tab:metodologia:entregables}
da el estado de cada entregable de la línea base y dónde está su
evidencia. Cinco se completan, el del catálogo con ocho APIs en lugar
de diez a doce. El de la comparación queda parcial y el de la memoria
sigue en curso. Se añade un
entregable que la línea base no preveía, la contribución a SPHERE.

\begingroup
\setlength{\tabcolsep}{3pt}
\setlength{\LTcapwidth}{\textwidth}
\begin{longtable}{|>{\footnotesize}p{0.09\textwidth}|>{\footnotesize}p{0.31\textwidth}|>{\footnotesize}p{0.10\textwidth}|>{\footnotesize}p{0.2\textwidth}|>{\footnotesize}p{0.21\textwidth}|}
\hline
\raggedright \textbf{Paquete} & \raggedright \textbf{Entregable} & \raggedright \textbf{Objetivo} & \raggedright \textbf{Evidencia} & \raggedright \textbf{Estado} \tabularnewline
\hline
\endfirsthead
\multicolumn{5}{l}{\footnotesize\textit{(continuación)}} \tabularnewline
\hline
\raggedright \textbf{Paquete} & \raggedright \textbf{Entregable} & \raggedright \textbf{Objetivo} & \raggedright \textbf{Evidencia} & \raggedright \textbf{Estado} \tabularnewline
\hline
\endhead
\raggedright PT0 & \raggedright Registro de decisiones y repositorio versionado & \raggedright --- & \raggedright Repositorio de datos & \raggedright Completado \tabularnewline
\hline
\raggedright PT1 & \raggedright Estado del arte & \raggedright --- & \raggedright Capítulo~\ref{cap:estadoDelArte} & \raggedright Completado \tabularnewline
\hline
\raggedright PT2 & \raggedright Criterio de inclusión, conjunto de APIs y Datasheet piloto de Postmark & \raggedright OBJ-01 & \raggedright Secciones~\ref{sec:propuesta:seleccion} y~\ref{sec:propuesta:piloto} & \raggedright Completado \tabularnewline
\hline
\raggedright PT3 & \raggedright Datasheets del catálogo, inventario de elementos no soportados y propuestas de extensión del esquema & \raggedright OBJ-02 y OBJ-03 & \raggedright Sección~\ref{sec:propuesta:catalogo}, Capítulo~\ref{cap:evolucion} y repositorio de datos & \raggedright Completado con ocho APIs de las diez a doce previstas \tabularnewline
\hline
\raggedright PT4 & \raggedright Datasheets ejecutadas en PRIME4API y relación entre límites y \textit{rates} & \raggedright OBJ-04 & \raggedright Sección~\ref{sec:propuesta:prime4api} y Capítulo~\ref{cap:limites} & \raggedright Completado dentro del ciclo de PT3 \tabularnewline
\hline
\raggedright PT5 & \raggedright Catálogo de preguntas de decisión y procedimiento de comparación & \raggedright OBJ-05 & \raggedright Secciones~\ref{sec:comparacion:catalogo} y~\ref{sec:comparacion:procedimiento} & \raggedright Parcial: las preguntas que exigen HARVEY4API no se ejecutan \tabularnewline
\hline
\raggedright --- & \raggedright Colección de las ocho APIs con su pricing en Pricing2YAML, una maqueta (prototipo web de cómo se vería la colección en SPHERE) y una propuesta de sección, en la ficha de cada API de la colección, para los términos de consumo que su pricing no recoge & \raggedright OBJ-06 & \raggedright Sección~\ref{sec:propuesta:sphere} & \raggedright Añadido el 15 de octubre y cargado en SPHERE el 27 de octubre, en una colección privada, sin publicar \tabularnewline
\hline
\raggedright PT6 & \raggedright Memoria y defensa & \raggedright --- & \raggedright Este documento & \raggedright En curso: depósito el 2 de noviembre y defensa del 9 al 13 \tabularnewline
\hline
\noalign{\vskip\abovecaptionskip}
\caption[Entregables de la línea base y su estado]{Entregables de la línea base y su estado. La fila sin paquete es un entregable añadido después de la línea base.}
\label{tab:metodologia:entregables}
\end{longtable}
\endgroup

\paragraph{Alcance.} Las decisiones de modelado se registran como
reglas en el libro de criterios y no cuentan como desviaciones. La
Tabla~\ref{tab:metodologia:desviaciones} recoge las decisiones que
cambian el alcance, con su fecha y el motivo que consta en el registro
de decisiones.

\begingroup
\setlength{\tabcolsep}{3pt}
\setlength{\LTcapwidth}{\textwidth}
\begin{longtable}{|>{\footnotesize}p{0.07\textwidth}|>{\footnotesize}p{0.25\textwidth}|>{\footnotesize}p{0.33\textwidth}|>{\footnotesize}p{0.28\textwidth}|}
\hline
\raggedright \textbf{Fecha} & \raggedright \textbf{Previsto} & \raggedright \textbf{Ejecutado} & \raggedright \textbf{Motivo registrado} \tabularnewline
\hline
\endfirsthead
\multicolumn{4}{l}{\footnotesize\textit{(continuación)}} \tabularnewline
\hline
\raggedright \textbf{Fecha} & \raggedright \textbf{Previsto} & \raggedright \textbf{Ejecutado} & \raggedright \textbf{Motivo registrado} \tabularnewline
\hline
\endhead
\raggedright 3 ago & \raggedright Una versión nueva del esquema (v0.8 o v1.x) & \raggedright Datasheets conformes a ACM2YAML v0.3. Los elementos no soportados se registran, y de ellos salen las propuestas de extensión (Capítulo~\ref{cap:evolucion}) & \raggedright Mantener las Datasheets consumibles por PRIME4API y SPHERE sin bifurcar el esquema \tabularnewline
\hline
\raggedright 5 ago & \raggedright Revisar con cuenta propia el panel del proveedor solo en algunas APIs & \raggedright Se revisa el panel con cuenta propia en las ocho APIs, y ninguna Datasheet se cierra sin esa revisión (Sección~\ref{sec:propuesta:ciclo}) & \raggedright Si unas APIs se revisan más que otras, las diferencias entre sus Datasheets reflejan el esfuerzo dedicado a cada una y no lo que publica cada proveedor \tabularnewline
\hline
\raggedright 1 sep & \raggedright Diez a doce APIs & \raggedright Ocho APIs, con un mínimo de seis, cerradas sin sustituciones (Sección~\ref{sec:propuesta:seleccion}) & \raggedright Se fija junto con el calendario de la replanificación \tabularnewline
\hline
\raggedright 1 sep & \raggedright Tres fuentes de selección: un trabajo previo, las APIs ya modeladas y una consulta a un modelo de lenguaje & \raggedright Se ejecutan la segunda y la tercera, y la tercera sustituye a la primera (Sección~\ref{sec:propuesta:seleccion}) & \raggedright La memoria del trabajo previo llegó el 14 de septiembre, con la selección ya hecha \tabularnewline
\hline
\raggedright 6 sep & \raggedright Las Datasheets, en anexos de la memoria & \raggedright En el repositorio de datos. La memoria recoge los hallazgos & \raggedright Los ficheros de datos se publican en un solo medio, el repositorio \tabularnewline
\hline
\raggedright 7 sep & \raggedright Ejecución sobre el despliegue público del ecosistema & \raggedright PRIME4API en local, con las preguntas que ejecuta su motor (Sección~\ref{sec:comparacion:catalogo}). Las preguntas que exigen HARVEY4API quedan sin ejecutar & \raggedright Sin acceso al despliegue público \tabularnewline
\hline
\raggedright 15 oct & \raggedright Sin colección en SPHERE. Pricing2YAML, fuera del alcance & \raggedright Los ocho pricings en Pricing2YAML y una maqueta de la colección con el histórico de cada pricing (Sección~\ref{sec:propuesta:sphere}) & \raggedright Ampliación: la colección lleva el catálogo a la plataforma de la línea de investigación \tabularnewline
\hline
\raggedright 23 oct & \raggedright Cinco objetivos específicos. La recogida de los términos de consumo, implícita en el catálogo & \raggedright OBJ-02 nombra la recogida sistemática de los términos de consumo, y OBJ-06 recoge la contribución a SPHERE (Sección~\ref{sec:introduccion:objetivos}) & \raggedright Reformulación de objetivos para que reflejen el alcance ejecutado \tabularnewline
\hline
\noalign{\vskip\abovecaptionskip}
\caption[Desviaciones de alcance]{Desviaciones de alcance, por fecha de decisión en 2026. Cada fila recoge lo previsto, lo ejecutado y el motivo que consta en el registro de decisiones.}
\label{tab:metodologia:desviaciones}
\end{longtable}
\endgroup

Dos decisiones reducen el alcance: el catálogo de ocho APIs y las
preguntas que exigen HARVEY4API, que quedan sin ejecutar. Dos lo
amplían: la contribución a SPHERE y la recogida sistemática de los
términos de consumo como parte de OBJ-02. Las otras cuatro cambian cómo se
ejecuta o se presenta lo previsto, no cuánto se hace.

Lo que queda fuera pasa a trabajo futuro. La primera tarea es ejecutar
las preguntas que exigen HARVEY4API. La segunda está en la maqueta: la
ficha de cada API tiene dos pestañas, la del pricing y la del resto de
los términos de consumo, y cada una guarda sus versiones por separado,
así que falta una línea temporal única para toda la ficha.

Las dos últimas tareas tratan de \textit{skills}, instrucciones que sigue
un agente basado en un modelo de lenguaje para leer los términos de
consumo. La tercera es una \textit{skill} que interprete la diferencia
entre dos versiones de esos términos, pricing incluido, como un cambio del
dominio, por ejemplo un descuento anual. La cuarta es publicar las
\textit{skills} en un mercado en línea (\textit{marketplace}) y
conectarlas con SPHERE.
